Currently working as postdoctoral researcher at
\href{https://www.univ-rouen.fr/university-of-rouen-normandy/}{University of Rouen},
vcastor's research goes from \href{https://application.wiley-vch.de/books/sample/3527307486_c01.pdf}{QTAIM}
and \href{https://www.mdpi.com/1420-3049/25/17/4028}{IQA} analysis to \href{https://en.wikipedia.org/wiki/Nuclear_magnetic_resonance}{NMR} and
\href{https://pubs.acs.org/doi/pdf/10.1021/cr990029p?ref=article_openPDF}{Conceptual DFT} calculations.

VCastor has worked in the groups
of \href{https://iquimica.unam.mx/academicos/dr-tomas-rocha-rinza/}{Tom\'as Rocha Rinza},
\href{https://www.lct.jussieu.fr/pagesperso/contrera/}{Julia Contreras-Garc\'ia},
\href{https://www.lab-cobra.fr/annuaire/joubert-laurent/}{Laurent Joubert} and
\href{https://www.lab-cobra.fr/annuaire/tognetti-vincent/}{Vincent Tognetti},
focusing on hydrogen bonds, topological analysis of the electron density
and electron localisation function, also on the development of
polarisability and dipole moment under the QTAIM partition.

She has written the modules for QTAIM dipole moment and polarisability
for \href{https://www.scm.com}{SCM}
(\href{https://www.scm.com/amsterdam-modeling-suite/adf/}{ADF},
\href{https://www.scm.com/amsterdam-modeling-suite/periodic-dft-band-quantum-espresso/}{BAND}
and \href{https://www.scm.com/amsterdam-modeling-suite/dftb-mopac/}{DFTB}),
as well as refactoring other parts of the QTAIM analysis.

Her PhD thesis was focused on the chemical reactivity from the prespective
of the topological analysis of the density and Conceptual DFT, taking into
account the solvent effects. You can read
her \href{https://github.com/vcastor/PhDManuscript/blob/main/tesis.pdf}{PhD manuscript}.
